% Material suplementar, compilado no fim do manuscrito para as referencias cruzadas
% funcionarem. Para a submissao, separar em arquivo proprio.
\clearpage
\appendix
\setcounter{table}{0}
\setcounter{figure}{0}
\setcounter{section}{0}
\renewcommand{\thetable}{S\arabic{table}}
\renewcommand{\thefigure}{S\arabic{figure}}
\renewcommand{\thesection}{S\arabic{section}}

\begin{center}
\large\bfseries Supplementary Material
\end{center}

\section{Computing environment}\label{sec:s-repro}

XGBoost is deterministic within one environment, but its accuracy on this task changed with
the library version and the number of fitting threads. With code, data and seed fixed, sMAPE
ranged from 6.80\% to 6.98\% across seven released versions and five thread counts. CatBoost
did not move by more than a tenth of a basis point under the same variation. We therefore
pinned the XGBoost version and fitted it on one thread; the reported 6.88\% comes from that
environment. On a second machine with a different processor architecture, the same versions
and one thread gave 6.83\% for XGBoost and 6.41\% for its temperature variant (6.51\% in
the pinned environment). This
difference is small against every comparison between models in the paper, but it changes
the count of significant horizons for the temperature gain of XGBoost (main text, Temperature).

Prophet derives its interval bounds from a finite sample of posterior draws, which we seeded
at every window. Without the seed the bounds moved by 34 to 53 deaths between two fits of
the same window. Across machines, 63 of the 103 Prophet point forecasts differed by up to 49
deaths, which left coverage at 0.772 and moved sMAPE from 4.701\% to 4.700\%.

TabPFN was run through a hosted inference API, so its model version is outside the pinned
environment. The same protocol gave 6.04\% with an earlier client release and 5.85\% with
the current one.

\section{Boosting configurations}

% gerado por paper_v1b/build_tables.py a partir de paper/tables/tab7_variantes.tex -- nao editar a mao
\begin{table}[htbp]
\centering
\small
\setlength{\tabcolsep}{4pt}
\caption{Feature engineering variants for the two gradient boosting models, under the protocol of Table~\ref{tab:desempenho}. The last two columns compare each CatBoost variant with the seasonal naive method.}
\label{tab:variantes}
\begin{tabular}{lrrrrcc}
\toprule
& \multicolumn{2}{c}{CatBoost} & \multicolumn{2}{c}{XGBoost}
& \multicolumn{2}{c}{CatBoost vs seasonal naive} \\
\cmidrule(lr){2-3} \cmidrule(lr){4-5} \cmidrule(lr){6-7}
Variant & sMAPE & MASE & sMAPE & MASE & 95\% CI of $\Delta$ & DM \\
\midrule
Lags only (as reported) & 6.59 & 1.213 & 6.88 & 1.263 & $[-0.06, +0.72]$ & 0/6 \\
+ rolling window statistics & 6.79 & 1.248 & 6.95 & 1.275 & $[+0.16, +0.90]$ & 0/6 \\
+ Fourier calendar terms & 6.27 & 1.150 & 6.61 & 1.214 & $[-0.37, +0.39]$ & 0/6 \\
Seasonal differencing & 6.20 & 1.142 & 6.35 & 1.166 & $[-0.37, +0.23]$ & 0/6 \\
Differencing + windows + calendar & 6.14 & 1.129 & 6.44 & 1.181 & $[-0.43, +0.14]$ & 0/6 \\
As above, direct multi-step & 6.09 & 1.126 & 6.39 & 1.183 & $[-0.45, +0.10]$ & 0/6 \\
\midrule
Seasonal naive (benchmark) & \multicolumn{4}{c}{6.27 \quad 1.148} & --- & --- \\
\bottomrule
\end{tabular}

\vspace{0.5em}
\begin{minipage}{\textwidth}
\footnotesize
$\Delta$, CatBoost minus seasonal naive sMAPE, with a paired bootstrap 95\% CI. DM, horizons with $p<0.05$ out of six. Direct multi-step fits one model per horizon.
\end{minipage}
\end{table}


% gerado por paper_v1b/build_tables.py a partir de paper/tables/revisao_tab_optuna.tex -- nao editar a mao
\begin{table}[htbp]
\centering
\small
\setlength{\tabcolsep}{4pt}
\caption{Bayesian hyperparameter search for the two gradient boosting models.}
\label{tab:optuna}
\begin{tabular}{lrrrrrrr}
\toprule
& \multicolumn{3}{c}{sMAPE (\%)} & \multicolumn{2}{c}{Gain (pp)} & & \\
\cmidrule(lr){2-4} \cmidrule(lr){5-6}
Model & Untuned & Training-only & Test-exposed & Training-only & Test-exposed & Trials & Hours \\
\midrule
CatBoost & 6.59 & 6.42 & 6.35 & $+0.17$ & $+0.24$ & 40 + 60 & 2.1 \\
XGBoost & 6.88 & 7.02 & 6.25 & $-0.13$ & $+0.63$ & 40 + 60 & 1.1 \\
\midrule
Seasonal naive & \multicolumn{3}{c}{6.27} & --- & --- & --- & --- \\
\bottomrule
\end{tabular}

\vspace{0.5em}
\begin{minipage}{\textwidth}
\footnotesize
Training-only: 40 trials with the search restricted to the first 60 months, so no test window enters it. Test-exposed: 60 trials allowed to score on the test windows, an upper bound that a real procedure cannot reach. Gain is positive when the search improved on the untuned configuration of Table~\ref{tab:desempenho}.
\end{minipage}
\end{table}


Table~\ref{tab:variantes} changes one element of the boosting feature set at a time. Rolling
window statistics made both models worse. Fourier calendar terms and forecasting the
seasonal difference helped, and the best combination with one model per horizon reached
6.09\% for CatBoost against 6.27\% for the seasonal naive method, with a paired interval
containing zero and no significant horizon. In Table~\ref{tab:optuna}, the search restricted
to the training period improved CatBoost by 0.17 points and made XGBoost 0.13 points worse.
With the test windows exposed, CatBoost reached 6.35\% and XGBoost 6.25\%, 0.02 points ahead
of the seasonal naive method; these values come from configurations selected on the test set
and bound what tuning could achieve.

The best feature engineering variant also lost accuracy under the 60-month sliding window:
0.50 points for CatBoost and 0.41 for XGBoost. Only the aggregate metrics of that run were
kept, so it has no paired interval.

\section{TabPFN}

TabPFN received the same twelve lags as the boosting models. Its sMAPE of 5.85\% was 0.74
points ahead of CatBoost and 1.03 ahead of XGBoost. Against the seasonal naive method its
advantage was 0.42 points (interval $[-0.698, -0.128]$), and against the drift variant 0.33
points ($[-0.622, -0.031]$). The Diebold-Mariano test was significant only at the one-month
horizon ($p=0.0012$ and $p=0.0044$), with $p$ above 0.35 from the third horizon onward, so
the advantage is confined to the first month. SARIMA was ahead of TabPFN in four of six
horizons.

\section{COVID-19 period}

% gerado por paper_v1b/build_tables.py a partir de paper/tables/tab8_covid.tex -- nao editar a mao
\begin{table}[htbp]
\centering
\small
\setlength{\tabcolsep}{5pt}
\caption{sMAPE of the three most accurate models with and without the COVID-19 period in the test set.}
\label{tab:covid}
\begin{tabular}{lrrrrcc}
\toprule
& & \multicolumn{3}{c}{sMAPE (\%)} & \multicolumn{2}{c}{Leading three} \\
\cmidrule(lr){3-5} \cmidrule(lr){6-7}
Test period & Windows & Prophet & SARIMA & TimesFM & Pairs separated & Worst DM \\
\midrule
Full test period & 103 & 4.70 & 4.80 & 4.83 & 0 of 3 & 0 of 6 \\
Excluding 2020-03 to 2021-12 & 76 & 3.60 & 3.74 & 3.66 & 0 of 3 & 0 of 6 \\
Excluding 2020-01 to 2022-12 & 62 & 3.47 & 3.41 & 3.59 & 0 of 3 & 0 of 6 \\
\bottomrule
\end{tabular}

\vspace{0.5em}
\begin{minipage}{\textwidth}
\footnotesize
Windows whose test months overlap the excluded range are dropped whole, so the pairing is kept; no model is refitted. Pairs separated: paired comparisons among the three whose interval excludes zero and whose DM test reaches $p<0.05$ in at least three horizons. Worst DM: the largest count of significant horizons in any pair.
\end{minipage}
\end{table}


\section{Temperature and the calendar}

% gerado por paper_v1b/build_tables.py a partir de paper/tables/revisao_tab_temp_diffcal.tex -- nao editar a mao
\begin{table}[htbp]
\centering
\small
\setlength{\tabcolsep}{5pt}
\caption{Temperature gain for the boosting models with and without explicit calendar terms and seasonal differencing.}
\label{tab:tempdiffcal}
\begin{tabular}{lrrr}
\toprule
Model and feature set & Without temp. & With temp. & Gain (pp) \\
\midrule
CatBoost, lags only & 6.59 & 6.33 & $+0.258$ \\
CatBoost, + calendar and seasonal difference & 6.13 & 6.05 & $+0.076$ \\
\midrule
XGBoost, lags only & 6.88 & 6.51 & $+0.368$ \\
XGBoost, + calendar and seasonal difference & 6.37 & 6.38 & $-0.010$ \\
\bottomrule
\end{tabular}

\vspace{0.5em}
\begin{minipage}{\textwidth}
\footnotesize
Gain in percentage points, positive when temperature reduces error. The lags-only rows reproduce Table~\ref{tab:temperatura}.
\end{minipage}
\end{table}


\begin{figure}[htbp]
\centering
% gerado por scripts/build_paper_assets.py -- nao editar a mao
\begin{tikzpicture}
\begin{axis}[
  width=0.78\textwidth, height=4.6cm,
  xlabel={Reduction in sMAPE from the temperature covariate (pp)},
  xmin=-0.12, xmax=0.62,
  % com valores pequenos o pgfplots gera tick automatico em notacao cientifica
  % (-5 \cdot 10^-2) e os rotulos colidem. Tick e formato explicitos.
  xtick={-0.1,0,0.1,0.2,0.3,0.4,0.5,0.6},
  scaled x ticks=false,
  xticklabel style={/pgf/number format/fixed, /pgf/number format/precision=1},
  ymin=0.4, ymax=4.6, ytick={4,3,2,1}, yticklabels={Prophet,SARIMA,CatBoost,XGBoost},
  axis lines=left, tick align=outside, tick pos=left,
]
\draw[black!45, dashed, line width=0.6pt]
  (axis cs:0,0.4) -- (axis cs:0,4.6);
\addplot[draw=cprophet, line width=2.0pt, mark=none] coordinates {(-0.0808,4) (0.0129,4)};
\addplot[only marks, mark=*, mark size=2.2pt, draw=cprophet, fill=cprophet] coordinates {(-0.0346,4)};
\addplot[draw=csarima, line width=2.0pt, mark=none] coordinates {(0.0340,3) (0.2573,3)};
\addplot[only marks, mark=*, mark size=2.2pt, draw=csarima, fill=csarima] coordinates {(0.1358,3)};
\addplot[draw=ccatboost, line width=2.0pt, mark=none] coordinates {(0.0728,2) (0.4519,2)};
\addplot[only marks, mark=*, mark size=2.2pt, draw=ccatboost, fill=ccatboost] coordinates {(0.2581,2)};
\addplot[draw=cxgboost, line width=2.0pt, mark=none] coordinates {(0.1898,1) (0.5476,1)};
\addplot[only marks, mark=*, mark size=2.2pt, draw=cxgboost, fill=cxgboost] coordinates {(0.3683,1)};
\end{axis}
\end{tikzpicture}
\caption{Reduction in sMAPE from adding monthly minimum temperature under the climatology
policy, with 95\% paired bootstrap intervals. Positive values favour the model with
temperature.}
\label{fig:temp}
\end{figure}

\section{Error by horizon and seasonal profile}

\begin{figure}[htbp]
\centering
% gerado por scripts/build_paper_assets.py -- nao editar a mao
\begin{tikzpicture}
\begin{axis}[
  width=0.80\textwidth, height=5.2cm,
  xlabel={Forecast horizon (months)}, ylabel={sMAPE (\%)},
  xmin=0.7, xmax=6.3, xtick={1,2,3,4,5,6},
  axis lines=left, tick align=outside, tick pos=left,
  legend style={at={(0.02,0.98)}, anchor=north west, draw=none, fill=none,
                 font=\scriptsize, legend columns=3, column sep=4pt},
  ymin=4, ymax=8.6, ytick={4,5,6,7,8},
  % Sem o ymin e o ytick explicitos, a escolha automatica punha a marca mais baixa em
  % 5 e os tres modelos lideres ficavam abaixo dela, sem nenhuma referencia de leitura,
  % com o ponto do SARIMA em h=1 encostado no eixo.
]
\addplot[draw=cprophet, mark=*, mark size=1.6pt, line width=0.9pt, mark options={draw=cprophet, fill=cprophet}] coordinates {(1,4.3538) (2,4.4520) (3,4.6456) (4,4.8209) (5,4.9294) (6,5.0041)};
\addplot[draw=csarima, mark=*, mark size=1.6pt, line width=0.9pt, mark options={draw=csarima, fill=csarima}] coordinates {(1,4.1625) (2,4.6971) (3,4.7652) (4,4.8096) (5,5.1129) (6,5.2266)};
\addplot[draw=ctimesfm, mark=*, mark size=1.6pt, line width=0.9pt, mark options={draw=ctimesfm, fill=ctimesfm}] coordinates {(1,4.2967) (2,4.6208) (3,4.8998) (4,5.1074) (5,4.9729) (6,5.0975)};
\addplot[draw=ccatboost, mark=*, mark size=1.6pt, line width=0.9pt, mark options={draw=ccatboost, fill=ccatboost}] coordinates {(1,5.7154) (2,6.3801) (3,6.7560) (4,6.9079) (5,6.8022) (6,6.9744)};
\addplot[draw=cxgboost, mark=*, mark size=1.6pt, line width=0.9pt, mark options={draw=cxgboost, fill=cxgboost}] coordinates {(1,6.1161) (2,6.8891) (3,7.1997) (4,7.0746) (5,7.0017) (6,7.0013)};
\legend{Prophet,SARIMA,TimesFM,CatBoost,XGBoost}
\end{axis}
\end{tikzpicture}
\caption{sMAPE by forecast horizon, $n=103$ windows per cell.}
\label{fig:horizonte}
\end{figure}

\begin{figure}[htbp]
\centering
% gerado por scripts/build_paper_assets.py -- nao editar a mao
\begin{tikzpicture}
\begin{axis}[
  width=0.80\textwidth, height=4.8cm,
  axis y line*=left, axis x line*=bottom,
  xmin=0.5, xmax=12.5, xtick={1,...,12}, xticklabels={J,F,M,A,M,J,J,A,S,O,N,D},
  ylabel={Mean deaths per month}, ylabel style={text=black},
  scaled y ticks=false, ytick={6500,7000,7500,8000},
  yticklabel style={text=black, /pgf/number format/fixed,
                     /pgf/number format/1000 sep={,}},
  tick align=outside,
  legend style={at={(0,1.03)}, anchor=south west, draw=none, fill=none,
                 font=\scriptsize, text=black},
]
\addplot[draw=black!80, mark=*, mark size=2.0pt, line width=1.1pt]
  coordinates {(1,6969.9) (2,6272.8) (3,6856.3) (4,6764.0) (5,7678.1) (6,7971.6) (7,8247.2) (8,7893.4) (9,7268.9) (10,7201.0) (11,6817.7) (12,7018.1)};
\addlegendentry{Deaths per month}
\end{axis}
\begin{axis}[
  width=0.80\textwidth, height=4.8cm,
  axis y line*=right, axis x line=none,
  xmin=0.5, xmax=12.5,
  ylabel={Mean minimum temperature ($^\circ$C)}, ylabel style={text=black},
  yticklabel style={text=black},
  tick align=outside,
  legend style={at={(1,1.03)}, anchor=south east, draw=none, fill=none,
                 font=\scriptsize, text=black},
]
\addplot[draw=ctemp, mark=square*, mark size=1.7pt, line width=0.9pt, dashed,
         mark options={draw=ctemp, fill=ctemp, solid}]
  coordinates {(1,19.72) (2,19.79) (3,19.16) (4,17.54) (5,15.07) (6,14.06) (7,13.44) (8,13.83) (9,15.76) (10,16.93) (11,17.33) (12,19.14)};
\addlegendentry{Minimum temperature}
\end{axis}
\end{tikzpicture}
\caption{Mean cardiovascular deaths by month of year (left axis) and mean monthly minimum
temperature at station A701 (right axis), 2010--2023.}
\label{fig:sazonal}
\end{figure}
